% Transcription — fonds Alexandre Grothendieck, shelfmark 62, batch 1 (pages 1-20).
% Established from the facsimile archives/batches/62/batch-01.pdf (a copy of
% 62.pdf fetched through the project relay, no cover sheet: PDF page k =
% archivists' page k), per the skill transcribe-grothendieck. The folder is
% ninety-six pages in five batches; this is the first.
%
% Pass: Opus 5.5 (claude-opus-5-5), 2026-09-26 — first pass, unchecked against the pages by a human.
%
% Pages skipped: 2 and 17, blank cream sheets (show-through only); 19, the
% upside-down leaf of a typescript (« - 20 - », « n° 276 », §2 « Anneaux
% m-adiques noethériens », Artin-Rees) with nothing in his hand — an unrelated
% typed verso, skipped by the settled rule. Page 1 is a tan cover inscribed
% top right « Courbes elliptiques »; it takes no \page marker and its words
% become the section heading, with a \note.
%
% Pages summarised: none. Pages 18 and 20 are scratch sheets (tan paper, the
% other sides of typed leaves) carrying isolated computations among doodles;
% the computations are given in page order, the doodles are not transcribed.
%
% His pagination: none visible on these leaves; the only numbers are the
% archivists' pencilled ones bottom left. Numbered runs of his own:
% « Proposition 1 » (p. 4), « Proposition 2 » (p. 5), « Proposition 3 »
% (p. 6), with unnumbered Propositions and Corollaires between them; recorded
% where they fall.
%
% What the batch is about. Pages 3-16 are a connected draft, headed
% « Modules des courbes elliptiques » (p. 3): for an elliptic curve C over a
% locally noetherian S with unit section s, the invertible sheaf L = O(D) of
% the unit section, L_{C/S} = s^*(L) (the Lie sheaf), E_n = f_*(L^{⊗n});
% Prop. 1 (ranks of E_n by Riemann-Roch, L^{⊗n} ample iff n >= 1, very ample
% iff n >= 3); the filtration E_0 = E_1 ⊂ E_2 ⊂ E_3 ⊂ ..., the graded pieces
% L_n = E_n/E_{n-1} ≅ L^{⊗n} for n >= 2 and the graded algebra; Prop. 3 on
% the maps L ⊗ L_n -> L_{n+1} and the characteristic (a two-corollary block
% cancelled); the embedding C -> P(E_3), C as a cubic divisor, the inflexion
% tangent at the unit section; bases 1, y, x of E_3 adapted to the flag and
% the Weierstrass-type equation x^2 + a y^3 + bxy + cy^2 + dx + ey + f = 0,
% the discriminant delta, and the Proposition: an elliptic curve with a
% « base spéciale » of E_3 is the same as sections a..f with delta a unit;
% the affine scheme M = Spec Z[A..F][delta^{-1}] with the action of the
% triangular group, representing curves with special base (pp. 10-12); a
% finer normalisation by a trivialisation of L (group of pairs
% (lambda, triangular with lambda^2, lambda^3), equation x^2 - y^3 + ... ,
% remarks on 12th/24th roots of delta, p. 14); in residue characteristic
% ≠ 2 the symmetry x -> -x reduces the equation to x^2 = y^3 + ay^2 + by + c,
% C as a double cover of a Brauer-Severi scheme ramified along the section
% at infinity and an étale divisor of degree 3 (pp. 15-16). Page 18: j and
% j - 1 units, the cases j ≡ 1, j ≡ 0 with Kummer extensions of group Z/3,
% Z/2, b^3/c^2 = 27 (j-1)/j and the cubic x^3 + 27 (j-1)/j (x+1) = 0.
% Page 20: the discriminant of a cubic as prod f'(x_i), for a = 0,
% j = -27c^2/Delta and « Delta = b^3 - 27c^2 ? ».
%
% Notation fixed once for the batch: underlined sheaves as \underline{L},
% \underline{O}_S, \underline{V}_{C/S}; L_{C/S} as \mathcal{L}; E_n as
% \mathcal{E}_n and the graded pieces as \mathcal{L}_n (he underlines them
% irregularly; kept as written where it is clear); the gothic M of p. 10
% as \mathfrak{M}; affine space as \mathbf{E}^6_{\mathbf{Z}}. Inside
% formulas an illegible index or exponent is written « ? » (said in a note
% on p. 8), since \ill{} cannot sit in math.
%
% Hardest pages: 4-7 (the proof sketches between the statements), 11-12 and
% 14 (fast connective prose), 16 (a struck, boxed passage and a ten-line
% slanted marginal note). Readings to check first: the exponent of x on
% p. 8 and of Y in BXY^2 on p. 11; « P^3_S » on p. 9; the sign in
% delta(...) ≠ 0 on p. 10; the second fraction « (1/(j-1))^3 » on p. 18.

\documentclass[11pt,a4paper]{article}
% Le préambule commun à toutes les transcriptions du fonds Grothendieck.
%
% Il tient en une idée : **l'incertitude se compose, elle ne se résout pas.**
% Une transcription qui lisse un mot douteux a détruit la seule chose pour
% laquelle on la faisait — un lecteur doit pouvoir savoir, à chaque endroit,
% ce qui était sur le feuillet et ce qui a été ajouté. Les six macros qui
% suivent sont donc l'appareil critique, et non de la décoration.
%
% Les mêmes commandes sont reconnues par `scripts/render.mjs`, qui produit la
% vue de lecture du site. Une macro ajoutée ici doit l'être aussi là-bas,
% faute de quoi le rendu échouera bruyamment — ce qui est voulu : un rendu
% silencieusement faux est pire qu'un rendu absent.

\ProvidesPackage{grothendieck}[2026/08/08 Transcriptions du fonds Grothendieck]

\usepackage{amsmath,amssymb,amsthm}
\usepackage{mathrsfs}
\usepackage{stmaryrd}
% Les diagrammes commutatifs sont partout dans ce fonds : c'est la langue
% même de la géométrie algébrique de Grothendieck, et une transcription qui
% les rendrait en prose ne transcrirait rien.
\usepackage{tikz-cd}
% Français par défaut — c'est la langue des feuillets — mais l'anglais est
% chargé aussi : la traduction et le résumé sont des documents anglais, et la
% typographie française (espaces avant les deux-points) y serait une faute.
% Ces documents basculent par \selectlanguage{english} en tête de corps.
\usepackage[english,french]{babel}
\usepackage{fontspec}
\usepackage{geometry}
\geometry{a4paper,margin=25mm,marginparwidth=28mm}
\usepackage{marginnote}
\usepackage{xcolor}
% ulem, not soul. `soul`'s \st and \ul are built on a character-by-character
% scan that XeLaTeX's font machinery breaks: the document fails with an
% undefined control sequence rather than degrading. ulem does the same two jobs
% and is Unicode-engine safe. `normalem` stops it hijacking \emph, which the
% transcriptions use for Grothendieck's own emphasis.
\usepackage[normalem]{ulem}
\usepackage{hyperref}
\hypersetup{colorlinks=true,linkcolor=black,urlcolor=black,pdfborder={0 0 0}}

% --- Métadonnées du lot -----------------------------------------------------
% Reprises telles quelles par le script de rendu, qui les affiche en tête de
% la vue de lecture. Elles fixent ce que le fichier prétend couvrir : sans
% elles, une transcription flotte sans point d'attache dans un fonds de seize
% mille feuillets.

\newcommand{\folder}[1]{\gdef\grothFolder{#1}}
\newcommand{\batch}[1]{\gdef\grothBatch{#1}}
\newcommand{\leaves}[2]{\gdef\grothFirst{#1}\gdef\grothLast{#2}}
\newcommand{\foldertitle}[1]{\gdef\grothFoldertitle{#1}}
\gdef\grothFolder{?}\gdef\grothBatch{?}\gdef\grothFirst{?}\gdef\grothLast{?}\gdef\grothFoldertitle{}

% --- Le résumé --------------------------------------------------------------
%
% Il ouvre la lecture modernisée, sous le titre « Résumé », et il est composé
% en petit. Ce n'est pas un ornement : c'est le seul passage du document qui
% ne soit pas des mathématiques, et un lecteur doit pouvoir le reconnaître
% avant de l'avoir lu — puis le sauter s'il connaît déjà le sujet, ou s'y
% arrêter s'il ne le connaît pas. Le filet qui le suit marque la reprise du
% corps.

\newenvironment{resume}
  {\small\setlength{\parskip}{0.45em}}
  {\par\medskip\hrule\medskip}

% --- Mots-clés ---------------------------------------------------------------
%
% En anglais, en clôture du résumé de la lecture modernisée : le vocabulaire
% moderne sous lequel chercher ce que les feuillets construisent. Le manifeste
% du site (`npm run manifest`) les extrait pour étiqueter la cote entière —
% cette ligne est la source unique des tags, il n'y a pas de fichier à côté.
\newcommand{\keywords}[1]{\par\smallskip\noindent
  {\footnotesize\textsc{Keywords} --- \textit{#1}}\par}

% --- L'appareil critique ----------------------------------------------------

% Le numéro de feuillet, en marge. C'est l'ancre qui permet de régler un
% désaccord sur une formule en regardant *un* feuillet plutôt qu'une chemise
% de six cents. Le site s'en sert aussi pour tourner les pages du fac-similé
% quand on fait défiler la transcription.
\newcommand{\leaf}[1]{%
  \marginnote{\footnotesize\color{gray!70}\textsf{#1}}%
  \label{leaf:#1}\ignorespaces}

% Illisible : on ne devine pas. Un mot inventé qui se lit comme les autres est
% le pire résultat possible.
\newcommand{\ill}{\textcolor{red!60!black}{[\,\ldots\,]}}

% Lecture douteuse : le mot est donné, et signalé comme incertain.
\newcommand{\uncertain}[1]{\uline{#1}}

% Ajout de l'éditeur — une lettre manquante, un mot sous-entendu.
\newcommand{\add}[1]{\textcolor{blue!50!black}{[#1]}}

% Biffé par Grothendieck lui-même. Ce qu'il a rayé fait partie du document :
% une rature dit souvent où la pensée a changé de direction.
\newcommand{\struck}[1]{\sout{#1}}

% Note du transcripteur, sur l'état matériel du feuillet.
\newcommand{\note}[1]{\par\smallskip{\small\color{gray!60!black}\textit{[#1]}}\par\smallskip}

% Note marginale de Grothendieck, distincte du corps du texte.
\newcommand{\marginal}[1]{\par\smallskip\noindent\hspace{1em}%
  {\small\color{gray!60!black}#1}\par\smallskip}

% --- Titre ------------------------------------------------------------------

\AtBeginDocument{%
  \begin{center}
    {\large\bfseries Fonds Alexandre Grothendieck --- cote n\textsuperscript{o}~\grothFolder}\\[2pt]
    {\itshape\grothFoldertitle}\\[4pt]
    {\small feuillets \grothFirst--\grothLast\ (lot \grothBatch)}\\[2pt]
    {\footnotesize\color{gray!60!black}Transcription établie d'après le fac-similé numérisé
     par l'Université de Montpellier.\\
     Les lectures incertaines sont soulignées, les illisibles notées [\,\ldots\,],
     les ajouts entre crochets.}
  \end{center}
  \vspace{1em}}

\endinput


\folder{62}
\batch{1}
\pages{1}{20}
\dating{[à partir de 1963]}
\watermark{Édition de démonstration}
\foldertitle{Courbes elliptiques (généralités - vieille rédaction) : notes manuscrites (s.d.)}

\begin{document}

\section*{Courbes elliptiques}

\note{inscrit de sa main, en haut à droite du feuillet de garde (page 1), qui ne
porte rien d'autre ; le feuillet ne reçoit pas de numéro de page}

\section*{Modules des courbes elliptiques}

\note{titre de sa main, souligné, en tête de la page 3}

\page{3}
Soit $C$ une « courbe elliptique » homogène sur $S$ loc. noethérien. La section
unité $s$ définit un bon diviseur $D$ \add{transversal aux fibres}\note{« transversal
aux fibres » est écrit au-dessus de la ligne et appelé par un trait} sur $C$,
\uncertain{d'où} un \struck{fibré} \add{faisceau inv.\ $\underline{L}$}, dont les
\add{germes de} sections sont des \struck{fonctions} germes de fonctions sur $C$
(précisément celles qui sont des « \uncertain{multiples} » de $D$). Notons que
\struck{$s^*(\underline{L}) \simeq$}, $C$ étant un schéma en groupes
\note{la formule inachevée $s^*(\underline{L})\simeq$ est barrée de hachures ; la
ligne suivante, « Proposition 1 Les $\underline{L}^n$ ($n\in\mathbf{Z}$) sont
\ill{} \ill{} [plats sur les \ill{}] », est biffée d'un trait, et « Proposition 1 »
est relié par une boucle au mot « dimension » plus bas, où la proposition est
reprise à la page 4}
\[
\underline{V}_{C/S} = f^*\bigl(s^*(\underline{V}_{C/S})\bigr) = f^*(\mathcal{L}_{C/S})
\qquad [\,\mathcal{L} = \text{faisceau de Lie}\,]
\]
\[
s^*(\underline{L}) \simeq s^*(\underline{L}\otimes_{C}\underline{O}_{S'})
\simeq s^*(\underline{N}_{C/S'}) \simeq s^*(\underline{V}_{C/S}) = \mathcal{L}
\]
\note{le $S'$ désigne manifestement la section unité vue comme sous-schéma de $C$ ;
le dernier signe, au bord de la feuille, se lit mal}
ce faisceau sera noté $\mathcal{L}_{C/S}$ [« Lie »]. On a
\[
s^*(\underline{L}^{\otimes n}) \simeq \mathcal{L}_{C/S}^{\otimes n}.
\]
On a une section \uncertain{canonique} de $\underline{L}$, savoir la section
constante $1$ de $\underline{O}_C$,
\[
\sigma\colon \underline{O}_C \longrightarrow \underline{L}, \qquad
\sigma\in\Gamma(C,\underline{L})
\]
qui définit des homomorphismes (\uncertain{d'ailleurs injectifs})
\[
\cdots \longrightarrow \underline{L}^{\otimes n} \longrightarrow
\underline{L}^{\otimes(n+1)} \longrightarrow \cdots \longrightarrow
\underline{L}^{\otimes(n+k)} \longrightarrow \cdots
\]
\marginal{$\subset$ faisceau des fonctions \uncertain{rationnelles} \ill{} sur $C/S$}
d'où \struck{\ill{}} des homomorphismes
\[
\cdots \longrightarrow f_*(\underline{L}^{\otimes n}) \longrightarrow
f_*(\underline{L}^{\otimes(n+1)}) \longrightarrow \cdots
\]
i.e.
\[
\cdots \longrightarrow \mathcal{E}_n \longrightarrow \mathcal{E}_{n+1} \longrightarrow \cdots
\]
en posant
\[
\mathcal{E}_n = f_*(\underline{L}^{\otimes n}).
\]

\page{4}
\emph{Proposition 1.} Les faisceaux $\underline{L}^{\otimes n}$ sont coh.\ plats en
toute dimension. Le faisceau $f_*(\underline{L}^{\otimes n})$ est localement libre,
de rang $0$ si $n<0$, $1$ si $n=0$ ou $1$, $n$ si $n\geqslant 1$. \struck{\ill{}} Le
faisceau $\underline{L}$ est ample, et de façon précise $\underline{L}^{\otimes n}$
est \struck{très} \emph{ample} \struck{pour} ssi $n\geqslant 1$, très ample ssi
$n\geqslant 3$.
\note{la proposition est encadrée d'un trait vertical à gauche ; « ssi » est écrit
au-dessus d'un « pour » biffé}

La formule de RR, où $g=1$, d'où $2g-2=0$, nous montre que \emph{tout faisceau}
\struck{loc.\ libre} inversible sur $C$, de degré \add{relatif}\note{« relatif » en
interligne} $d\neq 0$, est coh.\ plat, et que sa dimension est donc nulle si
$d<0$, et $d$ \struck{\ill{}} si $d\geqslant 1$. Donc la première assertion est
démontrée pour $n\neq 0$, et pour $n=0$ elle est connue par ailleurs, car
$\underline{L}^{\otimes 0}=\underline{O}_C$. \ill{} \ill{} \ill{} \ill{} faisceau
inv.\ sur $C$ de degré relatif \ill{} \ill{} ssi $d\geqslant 0$, donc
$\underline{L}^{\otimes n}$ \ill{} \ill{} ssi $n\geqslant 1$.
\struck{Et pour \ill{} \ill{} $\underline{L}$ \ill{} de degré relatif $d>0$, il
\ill{} \ill{}, \ill{} \ill{} \ill{} $\underline{L}^{\otimes n}$ est \ill{} pour
$n\geqslant 0$.}
Si $\underline{L}^{\otimes n}$ est très ample, on a $n\geqslant 1$, d'ailleurs on
ne peut avoir $n=1,2$ car alors $\mathbf{P}(f_*(\underline{L}^{\otimes n}))$ serait
de \uncertain{dimension} relative $0$ ou $1$, donc $C$ ne pourrait s'y immerger
(car il y serait isomorphe, absurde), donc il faut $n\geqslant 3$. \ill{} \ill{},
on \ill{} \ill{} \ill{} \ill{} $2$ \ill{} \ill{}, \ill{} \ill{} \ill{} \ill{} \ill{}
$S$-\uncertain{géom}. \ill{}, \ill{} \ill{} \uncertain{composés} \add{\ill{}}.
\struck{Mais alors D'ailleurs,} \struck{\ill{} $\underline{L}$ est \ill{} \ill{}
\ill{} \ill{} \ill{}}
Mais \struck{alors} comme $3 > (2g-2)+2 = 2g = 2$, un raisonnement \ill{} n'utilisant
que RR donne le résultat.
\marginal{prop.\ \uncertain{préliminaire} \ill{} \ill{} faisceau \ill{} \ill{}}
\note{note marginale en biais, à gauche, au crayon ; tout le passage depuis
« La formule de RR » est marqué d'un trait vertical}

\page{5}
\emph{Proposition.} On a une suite \ill{} \struck{\ill{}} d'injections
\[
\mathcal{E}_0 = \mathcal{E}_1 \subset \mathcal{E}_2 \subset \mathcal{E}_3 \subset \cdots
\]
les $\mathcal{E}_i$ sont des faisceaux loc.\ libres de rang $i$
pour $i\geqslant 1$. D'ailleurs
\[
\mathcal{E}_0 = \underline{O}_S .
\]

\emph{Proposition 2.} $\mathcal{E}_0=\mathcal{E}_1=\underline{O}_S$, d'autre part
pour $i\leqslant j$ l'injection
$\mathcal{E}_i \to \mathcal{E}_j$ est \struck{loc.}\ universellement injective.
\struck{(pour $n\geqslant 1$, \ill{}) \ill{} \ill{} \ill{}} \ill{} alors \ill{}
\ill{} \ill{} $\otimes$ \ill{} isomorphisme
\note{passage surchargé d'ajouts et de ratures en interligne, lu seulement par
fragments}
Soit \struck{\ill{} faisceau}
\[
\mathcal{L}_n \simeq \mathcal{E}_{n}/\mathcal{E}_{n-1}
\]
\struck{\ill{} isomorphe à $\mathcal{L}$}
\struck{et $\mathcal{E}_{n+1}/\mathcal{E}_n$ est isomorphe à $\mathcal{L}^{\otimes n}$}
\[
\mathcal{L}_n\otimes\mathcal{L}_m \simeq \mathcal{L}_{m+n} \quad \text{si } n,m
\underset{\neq 1}{\geqslant 0}
\]
\struck{\ill{} \ill{}} \add{correspondant} sur $\coprod_{n\geqslant 0}\mathcal{L}_n$
\struck{\ill{}} d'algèbre graduée comm.
\note{l'indice de $\mathcal{E}_{n}$ au numérateur est surchargé ; les lettres
$\mathcal{L}$ de cette page sont soulignées}

Pour la première assertion, \uncertain{because} la platitude coh.\ \ill{} \ill{}
\ill{} \ill{} d'une \ill{}, c'est trivial car \ill{} \ill{} une injection.
Notons que
\[
\mathcal{E} = \varinjlim \mathcal{E}_n = f_*(\varinjlim \underline{L}^{\otimes n})
\]
est d'ailleurs une $\underline{O}_S$-\emph{algèbre}, dont la connaissance
\uncertain{équivaut} à la connaissance de
$\underline{\Gamma}_*(f,\underline{L})=\coprod \mathcal{E}_n$\note{un signe illisible en exposant de $\Gamma$}, \ill{} on a
\[
\mathcal{E}_m\times\mathcal{E}_n \longrightarrow \mathcal{E}_{m+n}
\]
\ill{} précisément induisent la multiplication de $\mathcal{E}$. On a, \ill{}
\ill{} \ill{} \uncertain{gradué associé} à la filtration croissante donnée, une
algèbre \ill{} $\coprod_{n\geqslant 0}\underline{\mathcal{L}}_n$, où
$\underline{\mathcal{L}}_n \simeq \mathcal{E}_n/\mathcal{E}_{n-1}$ pour
$n\geqslant 0$, \ill{} \ill{} pour $n=0$ \ill{} $\underline{O}_S$. Il faut prouver
que
\[
\underline{\mathcal{L}}_n\otimes\underline{\mathcal{L}}_m \longrightarrow
\underline{\mathcal{L}}_{n+m}, \qquad n,m \underset{\neq 1}{\geqslant 0},
\]
\struck{\ill{}} est bien isom, ou encore surjectif, \struck{\ill{} \ill{}
\ill{}} \struck{\ill{} \ill{} \ill{} \ill{}} \ill{} \ill{} $n,m\geqslant 2$, d'où
\ill{} \ill{}.
\note{les deux dernières lignes de la page sont en partie biffées et se
chevauchent}

\page{6}
\ill{} \ill{} à \ill{}.

\emph{Corollaire.} Si $n\geqslant 2$, $\underline{\mathcal{L}}_n$ est isomorphe par
l'homomorphisme induit à $\underline{\mathcal{L}}_2^{\otimes n'}$ si $n$ pair, à
$\underline{\mathcal{L}}_2^{\otimes n'}\otimes\underline{\mathcal{L}}_3$ si $n$
impair (où $n'\geqslant 1$ dans le premier cas, $n'\geqslant 0$ dans le deuxième).

Mais on a plutôt intérêt à prouver directement ceci : on est ramené aussitôt par
platitude coh.\ au cas où $S=\operatorname{Spec}(k)$, on peut alors \ill{} \ill{}
\ill{} de \ill{} \ill{} \ill{}, \ill{} \ill{} de \ill{} \ill{} \ill{}, \ill{}
\ill{} \ill{} \ill{} \ill{} \ill{} \ill{} \ill{} \ill{}.

\emph{Corollaire.} $\mathcal{E}_3$ engendre $\Gamma_*(f,\underline{L})$.

Considérons les opérations de $\underline{V}$ sur les germes de fonctions
\uncertain{rationnelles} sur $C$ [\uncertain{rendre ceci défini}], on constate que
$\underline{V}$ \ill{} opérations \uncertain{induisent} des homomorphismes
\[
\underline{\mathcal{L}}\otimes\underline{\mathcal{L}}_n \longrightarrow
\underline{\mathcal{L}}_{n+1}, \qquad n\geqslant 2 .
\]

\emph{Proposition 3.} \struck{(Les)} \add{Soit $s\in S$, $p$ l'exposant car.\ de
$k(s)$.}\note{ajout en interligne, relié par un trait ; la lettre lue $s$ est
écrite comme un $\beta$} \add{Les} homomorphismes
$\underline{\mathcal{L}}\otimes\underline{\mathcal{L}}_n\to\underline{\mathcal{L}}_{n+1}$
\ill{} isomorphismes (\uncertain{voir} \struck{\ill{}} \ill{} \ill{}) \ill{}
\ill{} ssi $p$ est premier à $n$.
\note{l'énoncé est marqué d'un trait vertical à gauche ; « isomorphismes » est
souligné}

On est ramené au cas où \ill{} \ill{} $S=\operatorname{Spec}(k)$, et alors
c'est une vérification locale : si $D$ \ill{} une \struck{diff} dérivation
\ill{} \ill{} \ill{} \ill{} \ill{}, \add{$x\in C$} \ill{} si $f\in{}$\ill{} sur $C$
\ill{} \ill{} $x$ \ill{} \ill{} \ill{} \ill{},

\page{7}
\struck{\ill{} divisible par $p$} alors \struck{$Df$} $Df$ \ill{} \ill{} $x$
\ill{} pôle d'ordre \ill{} \ill{} \ill{}, et \ill{} : \ill{} ssi $p\nmid n$.

\struck{\emph{Corollaire 1}\quad Si $k(s)$ est de car.\ nulle, alors}
\[
\underline{\mathcal{L}}\otimes\underline{\mathcal{L}}_n \simeq
\underline{\mathcal{L}}_{n+1} \quad \text{pour } n\geqslant 2, \text{ donc}
\]
\struck{$\underline{\mathcal{L}}_3\otimes\mathcal{L}\otimes\underline{\mathcal{L}}_2$,
$\mathcal{L}_4$} \add{$\mathcal{L}_4\simeq\mathcal{L}_2\otimes\mathcal{L}^2$}
\[
\mathcal{L}_n \simeq \mathcal{L}^{n-2}\otimes\mathcal{L}_2, \qquad \text{donc pour }
n\geqslant 2,
\]
\struck{donc \ill{} $\mathcal{L}_2$ \ill{}}

\emph{Corollaire 2}\quad Si $k(s)$ est de car.\ $p\neq 2$, alors
\struck{$\mathcal{L}_3\simeq\mathcal{L}_2\otimes\mathcal{L}$}, 
$\underline{\mathcal{L}}_4\simeq\underline{\mathcal{L}}_3\otimes\mathcal{L}$,
$\mathcal{L}_6\simeq\underline{\mathcal{L}}_5\otimes\underline{\mathcal{L}}$ etc.
[$\underline{\mathcal{L}}_{2n}=\underline{\mathcal{L}}_{2n-1}\otimes\mathcal{L}$ si
$n\geqslant 2$], \ill{} \ill{}
\[
\mathcal{L}_2 \simeq \mathcal{L}^{\otimes 2} .
\]
\note{les deux corollaires sont encadrés ensemble, et le cadre est traversé par
deux longs traits courbes, dont l'un part du titre biffé « Corollaire 1 » : le
bloc paraît annulé en entier. Dans « $\mathcal{L}_2\simeq\mathcal{L}^{\otimes 2}$ »
l'indice est surchargé. Sous cette page, $\mathcal{L}$ sans indice est
$\mathcal{L}_{C/S}$ de la page 3, écrit tantôt souligné, tantôt non}

\emph{Corollaire.} Si $S$ est de car.\ $p>0$ [i.e.\ $p\,\underline{O}_S=0$] et
si $n$ est un multiple de la caractéristique, alors
$\mathcal{L}\otimes\underline{\mathcal{L}}_n\to\underline{\mathcal{L}}_{n+1}$ est
\emph{nul}.

\struck{Prop.}\ Considérons la suite exacte canonique
\[
\cdots\longrightarrow\underline{L}^{\otimes n}\longrightarrow
\underline{L}^{\otimes(n+1)}\longrightarrow
(\underline{V}_{C/S})^{\otimes(n+1)}\otimes_{\underline{O}_C}\underline{O}_{S'}
\longrightarrow 0
\]
on en conclut la suite exacte
\[
0\longrightarrow\mathcal{E}_n\longrightarrow\mathcal{E}_{n+1}\longrightarrow
\mathcal{L}^{\otimes(n+1)}\longrightarrow R^1f_*(\underline{L}^{\otimes n})
\]
\note{au-dessous du dernier terme, « $\| \; 0$ si $n\geqslant 1$ »}
d'où
\[
\boxed{\ \mathcal{L}_{n+1}\simeq\mathcal{L}^{\otimes(n+1)} \qquad n\geqslant 2\ }
\]
\marginal{intervertir avec la prop.\ 2}
et la multiplication dans $\mathcal{A}$ \ill{} \ill{}, \uncertain{envisagée}
dans prop.\ 2 \ill{} \ill{} par la multiplication \uncertain{évidente} \ill{}
\ill{} \ill{} des $\mathcal{L}_n$, $n\geqslant 2$
\note{dans l'encadré, l'indice $n+1$ de $\mathcal{L}_{n+1}$ est surchargé ; la
lettre après « dans » est lue $\mathcal{A}$ sous réserve}

\page{8}
Considérons maintenant l'immersion projective
\[
C \longrightarrow \mathbf{P}(\mathcal{E}_3)
\]
définie par $L^{\otimes 3}$ en vertu de prop.\ 1. Alors, $C$ devient un diviseur
positif sur $\mathbf{P}(\mathcal{E}_3)$ \add{de degré relatif 3,} simple sur $S$,
il est donc défini par une \add{section sur $S$} \struck{élément} de l'espace
\uncertain{projectif} des 3-formes
\[
\mathbf{P}_3(\mathbf{P}(\mathcal{E}_3)) \simeq
\mathbf{P}(\operatorname{Symm}_3(\check{\mathcal{E}}_3)) .
\]
\struck{De plus, comme} \struck{Supposons que l'on se soit donné une}
\struck{trivialisation}\note{ce dernier mot, en tête de ligne, est encadré et
barré d'une diagonale}
Considérons la section \add{$s$} de $C$ \add{\uncertain{et} la droite relative $D$}
\note{« et la droite relative $D$ » est écrit en interligne}
\ill{} $\mathbf{P}(\mathcal{E}_3)$ et la \struck{droite \ill{} \ill{} \ill{}
\ill{}} \ill{} de $\mathbf{P}(\mathcal{E}_3)$ définis comme
$\mathbf{P}(\mathcal{E}_3/\mathcal{E}_2)$ et
$\mathbf{P}(\mathcal{E}_3/\mathcal{E}_1)$, alors on constate que $s$ est une section
de $C$, et que $D$ est une « tangente d'inflexion » \ill{} \uncertain{en ce}
point.
\marginal{croquis : une courbe et sa tangente d'inflexion}

\ldots

Considérons une \struck{trivialisation} \add{base de $\mathcal{E}_3$ compatible avec
le} \add{et commençant par $1_S\in\Gamma(S,\underline{O}_S)$} drapeau
$0\subset\mathcal{E}_1\subset\mathcal{E}_2\subset\mathcal{E}_3$ de
$\mathcal{E}_3$ [elle \struck{\ill{} \ill{} \ill{} $X,Y,X$} i.e.\ une base
$1,Y,X$ « dépend de $(3+3-1)$ ($=5$) paramètres »], \struck{\ill{} \ill{}}
\ill{} multiplication près \ill{},
\note{les ajouts en interligne sont appelés par des accolades ; les trois
inclusions du drapeau sont soulignées}
\struck{la forme $\varphi$ (définie à \ill{} près) prend la forme [si on remplace
$Z$ par $1$]}
\[
a x^{?} + b y^3 + c xy + d y^2 + e x + f y + g = 0
\]
\note{le passage « la forme $\varphi$ \ldots\ $=0$ » est encadré et annulé par des
traits ondulés qui le traversent ; l'exposant de $x$ est surchargé (un $2$ sur
un $3$, ou l'inverse) et noté « ? » ; dans les formules de ce lot, « ? » marque
un indice ou un exposant qu'on ne peut lire. Sous le cadre, une ligne coupée par le bord : « $x^3\,y^2$
\ldots\ \ill{} \ill{} \ill{} $\mathcal{E}_1$ \ill{} de $\mathcal{E}_2$ dans
$\mathcal{E}_3$ »}
\marginal{$\begin{pmatrix} 1 & \alpha & \beta \\ & \mu & \gamma \\ & & \nu \end{pmatrix}$}
\note{matrice triangulaire en marge gauche, à hauteur du cadre : le changement
de base compatible avec le drapeau}

\page{9}
\struck{qui est aussi la relation (unique} \add{\struck{à multiple près}}
\struck{, \ill{} \ill{} vérifiée aussi bien a priori) vérifiée par}
\uncertain{Remarquons} que
\[
1,\ y,\ x,\ y^2,\ xy,\ y^3
\]
forment une base de $\mathcal{E}_6$ (adaptée à sa filtration), \ill{} aura
donc une équation
\[
\varphi(x,y) = \boxed{\,x^2 + a y^3 + b xy + c y^2 + d x + e y + f = 0\,}
\]
\marginal{$\Delta(a,b,c,d,e,f)\in\mathcal{L}^{\otimes 12}$}
\note{cette formule, entourée d'un ovale, est à gauche sous $\varphi(x,y)$ ; le
signe devant $a y^3$ est surchargé}
qui est d'ailleurs une équation valable dans $\mathcal{E}_6$ considéré comme
composant de $\Gamma_*(f,\underline{L}^{?}) = \coprod_{i\geqslant 0}\mathcal{E}_i$,
\add{à valeurs dans $\mathcal{E}_{?}$}
satisfaite par les sections $x$, $y$ de $\mathcal{E}_3$
\struck{$x_1, x_2, x_3$} \add{\ill{} \ill{} $2$, $2$, section $1$}
\note{ajout en interligne, lu par fragments}
\struck{Cette équation} Si on désigne \ill{} \add{par} \struck{\ill{}} \add{cette}
équation \add{écrite sous la forme} \ill{} \ill{}
$\mathcal{E}_{?}=\mathcal{E}_1\otimes\mathcal{E}_{?}$ \ill{} \ill{}
\[
\Phi(x,y,z) = x^2 z + a y^3 + b xyz + c y^2 z + d x z^2 + f z^3 = 0
\]
\note{ainsi sur la page : le terme en $yz^2$ ne figure pas ; les exposants de
$x^2z$ et de $dxz^2$ sont surchargés, et un $z$ après $ay^3$ est biffé}
on \emph{constate} alors que $\Phi$ est la forme \ill{} définit
\struck{\ill{}} $C$ dans $\mathbf{P}(\mathcal{E}_3)$ de degré $3$ qui définit
\ill{} $\simeq\mathbf{P}^{3}_S$\note{l'exposant $3$ est douteux ; on attendrait $2$}. En effet si $\Phi'$ est la forme qui
définit $C$, alors $\Phi$ est un \emph{multiple de $\Phi'$} et
\struck{\ill{} \ill{}} donc un multiple par quelque \ill{} \ill{} inversible,
puisque $\Phi$ est « \ill{} \ill{} » en tout pt de $S$ \ill{} \ill{}
\uncertain{coefficient} des $1$ \ill{}.
\marginal{d'ailleurs, et \uncertain{éventuellement} \ill{}}
\marginal{\uncertain{précision} \ill{} pt \ill{} intersection complète}
\note{deux notes en biais dans la marge gauche, la première séparée du texte
par un trait vertical}

\page{10}
Réciproquement, partons d'une équation \ill{} \ill{}, où les \add{coeff.}
\struck{\ill{}} sont des sections de $\underline{O}_S$. On sait qu'\add{une telle
équation} \struck{\ill{}} définit un diviseur relatif \add{positif} de
$\mathbf{P}^2_S/S$, soit $C$. Il reste à écrire que c'est là une \emph{courbe
lisse} sur $S$. On sait qu'il faut pour cela que le discriminant de la forme
\[
\delta(a,b,c,d,e,f)
\]
soit une \emph{unité} partout. Ici le discriminant s'explicite : \ldots

On a obtenu un résumé

\emph{Proposition.} La donnée d'une courbe elliptique sur $S$, \add{munie}
\struck{\ill{}} d'une \add{« base spéciale »} \struck{« base adaptée »} de
$\mathcal{E}_3(C/S)$, \struck{\ill{}} est équivalente à la donnée de sections
$a,b,c,d,e,f$ de $\underline{O}_S$ satisfaisant $\delta(a,b,c,d,e,f)\neq 0$ ;
cette correspondance est compatible avec \struck{\ill{}} l'image inverse.
\note{l'énoncé est marqué d'un trait vertical à gauche. le signe
devant $0$ dans « $\delta(\ldots)\neq 0$ » est lu $\neq$ sous réserve ; il tient
lieu de la condition, dite plus haut, que $\delta$ soit une unité}

Les termes de ceci : Considérons le schéma affine de type fini sur $\mathbf{Z}$
\[
\mathfrak{M} = \operatorname{Spec}\bigl(\mathbf{Z}[A,B,C,D,E,F][\delta^{-1}]\bigr)
\]
qui est un \ill{} \ill{} affine de $\operatorname{Spec}\mathbf{Z}[A,B,C,D,E,F]
\simeq \mathbf{E}^6_{\mathbf{Z}}$. \struck{Faisons opérer dans \ill{}}
\note{la lettre gothique est lue $\mathfrak{M}$ ; $\mathbf{E}^6_{\mathbf{Z}}$ est
l'espace affine de dimension $6$, écrit avec un $\mathbb{E}$ ajouré}

\page{11}
\struck{puis \ill{} le groupe $G$, de type fini sur $\mathbf{Z}$, qui \ill{} est
représenté \ill{} \ill{} par} $\mathbf{E}_{\mathbf{Z}}$\note{ou
$\underline{O}_{\mathbf{Z}}$ ; la lettre, en tête de ligne, est surchargée}.
Considérons le sous-schéma en groupes de $\mathrm{GL}(3)_{\mathbf{Z}}$ qui
correspond aux matrices de la forme
\[
\begin{pmatrix} 1 & \alpha & \gamma \\ 0 & \lambda & \beta \\ 0 & 0 & \mu \end{pmatrix}
\]
(avec $\lambda\mu$ une unité donc), \add{il opère sur $\mathfrak{M}$} de la façon
évidente qui consiste à interpréter $A,B,C,D,E,F$ comme coeff.\ d'une forme
\struck{$X^3+AXY^2$}
\[
X^2 + AY^3 + BXY^{2} + CY^2 + DX + EY + F
\]
\note{l'exposant $2$ de $Y$ dans $BXY^2$ est une lecture douteuse ; la page 9 a $bxy$}
et à faire sur $X, Y$ la substitution affine
\[
Y' = \lambda Y + \alpha, \qquad X' = \mu X + \beta Y + \gamma
\]
\note{la première lettre de chacune des deux formules est surchargée, un $Y$
écrit sur un $X$ semble-t-il}
[\ill{} \ill{} \ill{} \ill{} \ill{} \ill{} \ill{}, \ill{} \ill{} \ill{}
\ill{} le vectoriel $\mathbf{E}_3$, donc dans $\mathbf{P}_3(\mathbf{E}_3)$, et qui
laisse invariant le \uncertain{sous-}\ill{} \struck{\ill{}} \add{\ill{}} de
\add{\uncertain{les lettres}} \struck{\ill{}} \ill{} $6$ des formes qui \ill{}
de la forme dite], d'où les opérations de $G$ sur $\mathfrak{M}$.
Alors le foncteur \add{\uncertain{l'un des} \ill{} \ill{} \ill{} près} qui à tout
$S$ \uncertain{associe} \ill{} \ill{} des courbes elliptiques « \uncertain{standard} »
\ill{} \ill{} sur $S$, est \emph{représentable} par $\mathfrak{M}$.
\marginal{\struck{Prop}}
De plus, pour deux \ill{} \ill{}

\page{12}
$\mathfrak{M}(S)$ définissant des courbes elliptiques isomorphes (\ill{}
\ill{}) il faut que \struck{\ill{}} \ill{} se \ill{} en une opération de
$G(S)$, \ill{} $G$ \ill{} \ill{} \ill{} sur $S$. Enfin, pour \struck{\ill{}}
toute courbe elliptique $C$ sur $S$, il existe une extension \struck{fid.\ plate}
\add{\uncertain{finie} \ill{} de type fini} $S'$ de $S$ [savoir le fibré
principal associé à $\mathcal{E}_3$ \ill{} le drapeau qui \ill{} \ill{} sur
$\mathcal{E}_3$] telle que $C_{S'}$ « provienne » d'un élément de
$\mathfrak{M}(S')$.

Ainsi, \ill{} \ill{} \ill{} \ill{}, le foncteur « \ill{} \ill{} à isomorphisme
près des courbes elliptiques sur $S$ » est \struck{\ill{}}
\uncertain{pro-}représentable, où $\mathfrak{D}$ désigne la catégorie des
morphismes simples surjectifs.
\note{la lettre gothique, $\mathfrak{D}$ sous réserve, est écrite sur un mot
biffé}
Si ce n'était pas pour \ill{} \uncertain{d'ailleurs}, \ill{} \ill{} \ill{}
\ill{} \ill{} \ill{} \ill{} \ill{}, \emph{car} \ill{} \ill{} \ill{} un
\uncertain{raisonnement} \ill{} \ill{} \ill{} \ill{} \ill{} \ill{} \ill{}
\ill{} \ill{} \ill{} $\delta$ \ill{} $C$ \ill{} \ill{} \ill{} \ill{} \ill{}
\ill{} \ill{} \ill{} \ill{} \ldots).

\page{13}
Une normalisation \ill{} \ill{} plus \uncertain{précise} est obtenue ainsi : on se
donne une trivialisation de $\mathcal{L}$ par une section partout \ill{} \ill{}
$D$ de $\mathcal{L}$, [dérivation invariante], on choisit alors la base $1, x, y$
de $\mathcal{E}_3$ de façon que l'image de $x$ dans $\mathcal{L}_2$ soit
$i_2(D)$, celle de $x$ dans $\mathcal{L}_3$ soit $i_3(D)$.
\struck{\ill{}}
\note{ainsi sur la page : « $x$ » les deux fois ; la seconde est sans doute
à lire $y$}
[i.e.\ on prend comme structure auxiliaire la donnée d'un \add{\ill{}}
$\mathcal{L}$, d'un $\mathcal{E}$ loc.\ libre de rang $3$, \ill{} \ill{} \ill{}
\ill{} \ill{} \ill{} \ill{} \ill{} \ill{} isom avec
$\mathcal{L}^0, \mathcal{L}^2, \mathcal{L}^3$ respectivement.
Le groupe d'automorphismes de cette structure est \struck{\ill{}} \ill{} le
sous-groupe de $\mathbf{G}_m\times\mathrm{GL}(3)$, formé des \add{couples}
\struck{\ill{}} de la forme
\[
\left(\lambda,\ \begin{pmatrix} 1 & \alpha & \beta \\ 0 & \lambda^2 & \gamma \\ 0 & 0 & \lambda^3 \end{pmatrix}\right)
\qquad
\begin{array}{l}
\lambda\in\Gamma(\mathbf{G}_{m,S}/S)=\Gamma(S,\underline{O}_S^*)\\
\alpha,\beta,\gamma\in\Gamma(S,\underline{O}_S)
\end{array}
\]
\struck{Le \ill{} de la structure ainsi \ill{} à \ill{} \ill{} \ill{} \ill{}}
\add{isomorphe} \struck{\ill{}} au sous-groupe \struck{\ill{}} des
\[
\begin{pmatrix} 1 & \alpha & \beta \\ 0 & \lambda & \gamma \\ 0 & 0 & \mu \end{pmatrix}
\quad\text{pour}\quad
\lambda = \left(\frac{\mu}{\lambda}\right)^2, \quad
\mu = \left(\frac{\mu}{\lambda}\right)^3 \qquad\text{i.e.}\quad
\boxed{\mu^2 - \lambda^3 = 0}
\]
\note{les exposants de $\mu^2-\lambda^3$, dans le cadre, sont surchargés}
L'équation de $C$ pour une telle base \struck{prend alors la forme}
[caractérisée par la condition $\underline{x^2-y^3\in\mathcal{E}_5}$] prend la
forme
\[
x^2 - y^3 + bxy + cy^2 + dx + ey + f = 0
\]
[i.e.\ on peut \uncertain{prendre} $a=-1$ dans l'équation plus haut
\add{\uncertain{avec} $\delta(b,c,d,e,f)$ une unité}]
\note{l'ajout est entouré, sous la ligne}

\page{14}
Il ne semble pas qu'on puisse obtenir de normalisation locale sans
\struck{\ill{}} restriction \ill{} \ill{} qui fasse intervenir moins de
paramètres, [si on veut qu'elle soit \uncertain{indépendante} \ill{} \ill{}
constantes \struck{\ill{}} \ill{} \ill{} \ill{} \ill{}, on \ill{} \ill{}]. Si on
\ill{} \ill{} \ill{} \ill{} de base finis et plats, on peut prendre la
normalisation plus fine en exigeant que $\delta(b,c,d,e,f)=1$ \ill{} \ill{}
\add{(intrinsèquement)} \ill{} \ill{} \ill{} trivialisation de $\mathcal{L}$
\struck{\ill{}} \ill{} \ill{} \ill{} sections \ill{} \ill{} \ill{}
\add{puissance} $12$-ième (\ill{} \ill{} $24$-ième) \add{de $\delta^{-1}$,
\uncertain{où} $\delta$ est} \struck{\ill{}} discriminant. \ill{} \ill{}
\ill{} \ill{} des \struck{\ill{}} normalisations plus poussées en car.\
$\neq 2$, puis en car.\ $\neq 3$, puis en \ill{}.
\note{le $12$ est écrit sur un $11$ ; le passage entier, dense et vite écrit,
se lit par fragments}

\page{15}
Transformons l'équation
\[
x^2 = y^3 + axy + by^2 + cx + dy + e
\]
par $s$, \ill{} en remplaçant par $sx=-x$, $sy=y$, on trouve
\[
x^2 = y^3 - axy + by^2 - cx + dy + e
\]
et en comparant on trouve $a=c=0$, donc l'équation \ill{} \ill{} \ill{} \ill{},
la forme
\[
\boxed{\,x^2 = y^3 + ay^2 + by + c\,}
\]
\note{dans l'encadré, les lettres $a$ et $b$ sont écrites sur d'autres ; $s$
est manifestement la symétrie $x\mapsto -x$, dont l'introduction n'est pas sur
cette page}
avec comme condition $\delta(a,b,c)\neq 0$, qui ici se \uncertain{traduit} par :
les \uncertain{solutions} de $P(y)=y^3+ay^2+by+c$ sont distinctes en tout pt de $S$.

Géométriquement, \struck{\ill{}} l'équation \uncertain{s'interprète} ainsi : la
\ill{} $C/\mathbf{Z}_2$ est un schéma de Brauer--Severi
\add{$P$}\note{le $P$ est écrit au-dessus de la barre de « Brauer--Severi »}
sur $S$, muni d'une section (correspondant à la section unité de $C$ sur $S$),
donc \struck{\ill{} et Brauer-\ill{}} \add{correspond} à un espace
\struck{\ill{}} \ill{} principal homogène sous un \struck{\ill{} \ill{} \ill{}}
fibré vectoriel \struck{\ill{}} \add{inversible} [il correspond à
$\underline{L}^{\otimes 2}$]. $C$ est un revêtement \struck{\ill{}} double
\uncertain{ramifié} \ill{} un diviseur \add{de $P/S$} étale sur $S$, qui se
décompose en la section à l'infini, et en un diviseur \add{de $P/S$} étale sur
$S$, \struck{\ill{} \ill{} \ill{}} de
$\underline{V}(\underline{L}^{\otimes 2})/S$ étale sur $S$, de degré relatif $3$.
\marginal{\ill{}}
\note{un mot en biais dans la marge gauche, en face du paragraphe, que borde un
trait vertical}

\page{16}
\struck{Supposons \ill{} \ill{} les corps résiduels de}
Notons que $\uncertain{-1}$\note{le symbole après « Notons que » est peu net}
\ill{} de $C$ sur $S$ \ill{} \ill{} \ill{} \ill{} \ill{} \ill{}, \ill{}
$C$ est un schéma abélien sur $S$. Cet \ill{} \ill{}
\add{est \uncertain{d'ordre} $2$, il} opère sur les $\mathcal{E}_i$ [de façon
compatible avec les structures, \add{et les opérations de $\mathcal{L}$ sur
$\mathcal{E}_\infty$} multiplication] et induit \struck{\ill{}} sur $\mathcal{L}$
le changement de signe.
\struck{\ill{} \ill{} \ill{} \ill{} \ill{} \ill{} \ill{} \ill{} les
caractéristiques résiduelles de \ill{} \ill{} $\mathcal{E}_1$, de $\mathcal{E}_2$
\ill{} il y aura \ill{} \ill{} \ill{} \ill{} \ill{} \ill{} \ill{}}
\note{deux lignes et demie barrées et entourées d'un long trait}
Supposons les car.\ résiduelles de $S$ $\neq 2$. Alors
\struck{1) Il existe un supplémentaire unique \ill{} \ill{} de $\mathcal{E}_1$
dans $\mathcal{E}_2$ \ill{} \ill{} qui \ill{} \ill{} \ill{} $s$ \ill{}}
\ill{} \ill{} \ill{} \ill{} \ill{} \ill{} $s$ qui \uncertain{trivialise} dans
$\mathcal{E}_2$ [\ill{} \ill{} \ill{} \ill{} \ill{}], et \ill{} restriction \ill{}
\ill{} \ill{} \ill{}, \add{$L_3$ \ill{}} \ill{} \ill{} $\mathcal{E}_2$ dans
$\mathcal{E}_3$ \add{\ill{}} il existe un supplémentaire \ill{} \ill{} $s$
\ill{} \ill{} \uncertain{symétrie}, \struck{\ill{}} qui \ill{} \ill{} \ill{} $L_3$.
\ill{} \ill{} : $\mathcal{L}.\,\mathcal{E}_2$, \ill{} \ill{} \ill{} $L_3$.
\struck{Alors $L_3^{\otimes 2}$ \ill{} un supplémentaire \ill{} de $\mathcal{E}_5$
dans $\mathcal{E}_6$, \ill{} $L_6$. \ill{} \ill{}}
\note{cette ligne et la suivante sont barrées et encadrées ; le mot « À
\uncertain{compléter} » est écrit dans le cadre}
\marginal{N.B.\ la \uncertain{condition} de $\mathbf{P}(\mathcal{E}_2)$ \ill{}
\ill{} \ill{} \ill{} \ill{} \ill{} \ill{} \ill{} \ill{} \ill{} \ill{} \ill{}
$\mathcal{L}^{\otimes 3}\otimes\mathcal{E}_3$ \ill{} \ill{} \ill{} \ill{}}
\note{note en biais dans la marge gauche, sur une dizaine de lignes, lue par
fragments ; un trait vertical la sépare du texte}
\ill{} \ill{} \ill{} \ill{} \ill{} \ill{} \ill{} \ill{} \ill{} \ill{} les
bases \uncertain{admissibles} de $\mathcal{E}_3$, \ill{} exigeant que $x$ soit
\struck{\ill{}} \ill{} de $\mathcal{E}_3$ ; donc le groupe des \ill{} \ill{}
base \ill{} est \ill{} \ill{} la forme
\[
\begin{pmatrix} 1 & \alpha & 0 \\ 0 & \lambda^2 & 0 \\ 0 & 0 & \lambda^3 \end{pmatrix}
\quad\text{i.e.}\quad
\begin{pmatrix} 1 & \alpha & 0 \\ 0 & u & 0 \\ 0 & 0 & v \end{pmatrix}
\quad\text{avec}\quad u^3 - v^2 = 0 .
\]
\note{dans la seconde matrice les deux coefficients diagonaux sont surchargés ;
$u$ et $v$ sont lus d'après l'équation qui suit}
\ill{} \ill{} \ill{} \ill{} \ill{} \ill{} \ill{} \emph{\ill{} \ill{} \ill{}
\ill{} \ill{}}.

Quelle est la normalisation correspondante de l'équation de $C$ ?

\page{18}
\note{feuille de brouillon, à l'encre, sans rédaction suivie : calculs isolés
répartis sur la page, séparés par des traits horizontaux, parmi des dessins
qui ne sont pas des mathématiques et ne sont pas transcrits. On les donne dans
l'ordre de la page, de haut en bas}

\struck{\ill{}} b

\emph{$j$ et $j-1$ unités}
\note{souligné deux fois}

\marginal{Cas $A$ \ill{} \emph{\ill{}}}
\begin{itemize}
\item $j$ non unité, \uncertain{puis} $j-1$, (\ill{} \ill{})
\quad $(1-j^{-1})^2\in A^6A^*$
\item $j-1$ non unité, $j$ \ill{} (\ill{} \ill{})
\quad $\bigl(\tfrac{\uncertain{1}}{j-1}\bigr)^3\in A^6A^*$
\end{itemize}
\note{les deux lignes sont réunies par une accolade ; l'exposant « $A^*$ » se lit
en indice au bord droit}

\begin{itemize}
\item $j\equiv 1$ \quad $b^3=0$, $c$ inversible \quad // si $A$ \ill{} \ill{}
\ill{} \ill{}, \ill{} \ill{} $c\in A^*/A^{*(3)}$
\item $j\equiv 0$ \quad $c^2=0$, $b$ inversible \quad // si $A$ \ill{} \ill{}
\ill{} \ill{}, \ill{} \ill{} $b\in A^*/A^{*(2)}$
\end{itemize}
\marginal{Ext.\ Kummer de gr $\mathbf{Z}/3$ \quad (flèche vers $j\equiv 1$)}
\marginal{Ext.\ Kummer de gr $\mathbf{Z}/2$ \quad (flèche vers $j\equiv 0$)}
\marginal{à vérifier / \uncertain{compléter}}

\[
j\equiv 1 \quad \Bigl(\frac{t}{t-1}\Bigr)^3 \in O^6O^* \qquad
\Bigl(\frac{t-1}{t}\Bigr)^2 \in O^6O^* \qquad
\sqrt[6]{\frac{t}{t-1}}
\]
\note{au-dessus, isolé à droite, « 1– » ; « \ill{} \ill{} » souligné à droite de la première formule ; un signe entre
les deux formules, qu'on ne peut lire}

\[
\sqrt[2]{\frac{t}{t-1}} \qquad \sqrt[3]{\frac{t-1}{t}} \qquad \tfrac12 \qquad \tfrac13
\]
\[
\lambda^{3/6}\cdot\lambda^{-2/6} = \lambda^{1/6} \qquad
s^6 = \frac{t}{t-1}
\]
\[
\frac{-27c^2}{b^3-27c^2} = j \qquad 27c^2(j-1) = b^3 j \qquad
\frac{b^3}{c^2} = j
\]
\[
\begin{array}{ccc} L & L^2 & L^3 \\ & b & c \end{array}
\qquad
\boxed{\ \frac{b^3}{c^2} = 27\,\frac{j-1}{j}\ }
\]
\[
b = 27\,\frac{j-1}{j} \qquad c = 27\,\frac{j-1}{j}
\]
\[
\boxed{\ x^3 + 27\,\frac{j-1}{j}\,(x+1) = 0\ }
\]
\note{dans « $b=27\frac{j-1}{j}$ » l'exposant de $b$ est une tache ; dans
« $27c^2(j-1)$ » le $1$ est surchargé}

\struck{$\underline{\mathrm{Hom}}(E,\ \ )$} \quad $f_*(\check{E})$ \qquad
$\underline{\mathrm{Hom}}_{\underline{O}}(E,\underline{O})$
\[
\underline{\mathrm{Hom}}_{\underline{A}}(E,\underline{A})\times
\underline{\mathrm{Hom}}_{\underline{O}}(\underline{A},\underline{O})
\]
\note{une flèche monte de cette dernière formule vers
$\underline{\mathrm{Hom}}_{\underline{O}}(E,\underline{O})$}

\page{20}
\note{feuille de calcul à l'encre, sur le même papier que la page 18 : le
discriminant du polynôme cubique}

\struck{$(x_1-x_2)(x_3$}
\[
\prod_{i\neq j}(x_i-x_j) = \prod_i \Bigl(\frac{f(x)}{x-x_i}\Bigr)_{x\to x_i}
= \prod_i f'(x_i)
\]
\note{dans la fraction, le numérateur $f(x)$ et le $x$ du dénominateur sont
surchargés}
\marginal{$(y+u)^3 + a(y+u)^2$ \quad $3u+a$ \quad $u=-\frac{a}{3}$}
\note{calcul en haut à droite, isolé par un trait courbe ; le premier $y$ est
écrit sur une autre lettre}
\[
f(x_i) =
\]
\[
3x^2 + 2ax + b
\]
\[
(3x_1^2+2ax_1+b)(3x_2^2+2ax_2+b)(3x_3^2+2ax_3+b)
\]

Supposons $\underline{\underline{a=0}}$ \ill{} \ill{}
\[
j = -\frac{27c^2}{\Delta}
\]
\marginal{$x_1x_2 \quad x_2x_3 \quad x_3x_1$ \quad // \quad $x_1x_2 \quad x_2x_3$}
\[
\Delta = \prod(3x_i^2+b) = 27(x_1x_2x_3)^2 + 9b\bigl(\textstyle\sum(x_1x_2)^2
+ 3b^2\sum x_i^2\bigr) + b^3
\]
\struck{$= -27c^2 + 9b(b^2 - 2\sum x_1^2x_2x_3$}
\[
(x_1x_2x_3)^2 = c^2
\]
\[
\textstyle\sum(x_1x_2)^2 = \bigl(\sum x_1x_2\bigr)^2 - 2\sum x_1^2x_2x_3
\]
\[
\textstyle\sum x_i^2x_2x_3 = \bigl(\sum x_i^2\bigr)\bigl(\sum x_2x_3
\]
\[
\boxed{\ \Delta = b^3 - 27c^2\ }\ ?
\qquad x_1^2x_2^2 + \qquad \textstyle\prod 3x_i^2+b
\]
\struck{$\bigl(\sum x_i^2\bigr)^2$}
\[
3\textstyle\sum(x_1x_2)^2 + \bigl(\sum x_1x_2\bigr)\bigl(\sum x_i^2\bigr)
\]
\struck{\ill{}} \struck{$x_1x_2$}
\note{ainsi sur la page : la parenthèse de $3b^2\sum x_i^2$ est fermée par une
accolade, celle de $(\sum x_2x_3$ ne l'est pas ; dans le cadre, un chiffre
devant $b^3$ est biffé, et le point d'interrogation est de sa main}

\end{document}
